% !TeX encoding = UTF-8
\chapter{Discussion and Conclusion}\label{chap:discussion}

\paragraph{Shortcut features and domain shift.} The Decision Tree separates attacks from benign flows almost perfectly on the CIC-IDS2017 test split (attack F1 0.986, false-positive rate 0.037\,\%). On our own benign traffic, however, the false-positive rate is 0.55\,\%, about 15 times higher. The false alarms come from the shortcut features of Section~\ref{sec:t3-interp}, \texttt{Destination Port} and \texttt{Init\_Win\_bytes\_*}, applied to automated plain-HTTP bursts that resemble the Selenium-driven attacks. The model partly learned ``automated HTTP to port 80 from a small-window server'', not the attack itself. Other differences add to the shift: the training flows come from an older CICFlowMeter version and a 2017 testbed with different hosts and a different application mix, whereas our capture contains QUIC and Steam traffic. The data itself also has quality problems, such as the broken label encoding, the duplicated column, infinite values, about 6,000 duplicate rows and two TCP flag columns that are identical. They match issues reported in the literature~\cite{engelen2021troubleshooting}, and each of them needed its own cleaning step in Chapter~\ref{chap:task3_preprocessing}.

\paragraph{Accuracy paradox.} Multiclass Logistic Regression reaches 98.54\,\% accuracy while detecting no attack, because 98.7\,\% of the flows are benign. Class weighting does not solve the imbalance either. It raises LR's attack recall from 0.85 to 0.99 but multiplies the false alarms by ten, and it lowers the multiclass macro-F1 of the tree.

\paragraph{Limitations.} We use a single train/test split without cross-validation. The tree is unpruned, and its deep, fine splits show up in the false-positive analysis. SQL injection has only 6 test flows, so its metrics are statistically weak. Brute Force and XSS cannot be separated well from flow statistics alone. Our own capture contains no attacks, so only the false-positive rate can be measured on our network. Dropping rows with infinite values slightly changes the benign distribution (0.08\,\%), and min-max scaling is sensitive to outliers. Finally, the real-time part replays recorded flows and is not live capture.

\paragraph{Possible improvements (not implemented).} Pruning the tree (\texttt{max\_depth}, \texttt{min\_samples\_leaf}), dropping the testbed-specific features (\texttt{Destination Port}, \texttt{Init\_Win\_bytes\_*}), and adding benign traffic from the deployment network to the training data should reduce the false positives.

\paragraph{Conclusion.} We turned the Thursday web-attack file of CIC-IDS2017 into a clean set of 164,179 flows with 60 features, selected 20 features per target, and trained Logistic Regression and Decision Tree classifiers. The Decision Tree outperforms Logistic Regression on both targets and detects all attacks in the binary test split, while accuracy hides that default multiclass LR detects none. On our own benign traffic, the tree's false-positive rate rose from 0.037\,\% to 0.55\,\%, so its benchmark score does not carry over to our network.
